\documentclass{article}

\PassOptionsToPackage{numbers}{natbib}
\usepackage[preprint]{neurips_2026}

\usepackage[utf8]{inputenc}
\usepackage[T1]{fontenc}
\usepackage{hyperref}
\usepackage{url}
\usepackage{booktabs}
\usepackage{amsfonts}
\usepackage{amsmath}
\usepackage{amssymb}
\usepackage{nicefrac}
\usepackage{microtype}
\usepackage{xcolor}
\usepackage{graphicx}
\usepackage{enumitem}
\usepackage{listings}
\usepackage{subcaption}
\usepackage{colortbl}
\usepackage[framemethod=default]{mdframed}

\lstset{basicstyle=\small\ttfamily, breaklines=true, frame=single, xleftmargin=1em}

\definecolor{bestgreen}{HTML}{D4EDDA}
\definecolor{warnyellow}{HTML}{FFF3CD}
\definecolor{badred}{HTML}{F8D7DA}
\definecolor{neutralblue}{HTML}{CCE5FF}

\title{Physical Feasibility vs.\ Prediction Accuracy:\\
  Why PQ Residuals Are the Primary Metric\\
  for AC Power-Flow Surrogates in Congestion Management}

\author{Changhun Kim}

\begin{document}
\maketitle

\begin{abstract}
  Voltage RMSE is the most common accuracy metric for GNN AC power-flow
  surrogates, but it is a proxy.  For congestion management and optimal
  power flow (OPF), the fundamental requirement is that the predicted
  voltage state satisfies the nodal power-balance equations.  This note
  formalises the distinction, establishes a metric priority order grounded
  in how NR-family solvers define feasibility, and applies it to a set of
  experimental runs on the LVN Heo1 grid.  The conclusion is that physical
  residuals ($\|\Delta P\|_\infty$, $\|\Delta Q\|_\infty$, mean/p95) are
  the primary metrics, and voltage RMSE is secondary.
\end{abstract}

%==============================================================================
\section{The AC Power-Flow Feasibility Problem}
%==============================================================================

An AC power-flow solution is a voltage vector $(V, \theta)$ that satisfies
the nodal power-balance equations at every constrained bus:

\begin{equation}
  \Delta P_i = P_i^{\rm set} - \mathrm{Re}\{V_i \overline{(YV)_i}\} = 0
  \quad \forall\, i \in \text{PV} \cup \text{PQ}
  \label{eq:dp}
\end{equation}
\begin{equation}
  \Delta Q_i = Q_i^{\rm set} - \mathrm{Im}\{V_i \overline{(YV)_i}\} = 0
  \quad \forall\, i \in \text{PQ}
  \label{eq:dq}
\end{equation}

This is the definition used by pandapower, PYPOWER, and MATPOWER~\cite{matpower}.
Feasibility is a set of \emph{equality constraints}, not a closeness measure.
Newton--Raphson declares convergence when the worst-case violation falls below
a tolerance:

\begin{equation}
  \|[\Delta P,\, \Delta Q]\|_\infty < \text{tol}
  \label{eq:nrconv}
\end{equation}

with slack-bus $P,Q$ and PV-bus $Q$ \emph{excluded} (they are free variables).
The default pandapower tolerance is $10^{-8}$\,MVA.

A GNN surrogate predicting $(\hat{V}, \hat{\theta})$ is \emph{physically
feasible} if and only if it satisfies \eqref{eq:dp}--\eqref{eq:dq} at the
specified setpoints.  Voltage RMSE does not test this.

%==============================================================================
\section{Why Voltage RMSE Is Insufficient for Congestion Management}
%==============================================================================

Congestion management and OPF require a predicted state from which
operationally meaningful quantities can be derived:

\begin{itemize}[noitemsep]
  \item Branch complex power flows $S_{ij} = V_i I_{ij}^*$
  \item Line current magnitudes $|I_{ij}|$ and thermal loading
  \item Voltage constraint violations ($V_{\min} \le |V_i| \le V_{\max}$)
  \item Feasibility margins and available transfer capacity
\end{itemize}

All of these are \textbf{functions of the voltage state}.  If the predicted
state does not satisfy the AC power-balance equations, these derived
quantities are unreliable even when voltage RMSE looks small.

The reason is nonlinearity and network stiffness.  In a multi-voltage LVN
grid, diagonal admittance entries can reach $|Y_{ii}| \sim 10^6$\,p.u.
A small voltage error $\epsilon$ at such a bus produces a mismatch:

\begin{equation}
  |\Delta S_i| \approx |Y_{ii}| \cdot |\epsilon| \sim 10^6 \cdot \epsilon
  \label{eq:stiff}
\end{equation}

So $\epsilon = 10^{-3}$\,p.u.\ in voltage can yield $|\Delta S_i| \sim
10^3$\,p.u.\ in power mismatch.  Voltage RMSE $= 10^{-3}$\,p.u.\ can
coexist with $\|\Delta P\|_\infty = 10^3$\,p.u. --- three orders of
magnitude apart.

This is directly confirmed by the LVN experiments: mean $|\Delta P|
\sim 4 \times 10^{-4}$\,p.u.\ for individual buses while
$\|\Delta P\|_\infty \sim 10^5$\,p.u.\ at the stiffest bus.

%==============================================================================
\section{Metric Priority Order}
%==============================================================================

\begin{table}[h]
\centering
\caption{Metric priority order for GNN AC-PF surrogates used in congestion
  management and OPF.  All mismatch quantities use the masked convention:
  slack excluded from $\Delta P$; slack + PV excluded from $\Delta Q$.}
\label{tab:priority}
\begin{tabular}{clp{8cm}}
\toprule
Priority & Metric & Rationale \\
\midrule
1 & $\|\Delta P\|_\infty$, $\|\Delta Q\|_\infty$ &
  Worst-case power-balance violation.  A single bus with large residual
  means the state is not AC-feasible.  This is the direct analogue of
  NR convergence criterion \eqref{eq:nrconv}. \\[4pt]
2 & mean / p95 $|\Delta P|$, $|\Delta Q|$ &
  Distribution of residuals across buses.  Mean shows typical quality;
  p95 shows near-worst quality without being dominated by one outlier.
  Together with Priority 1, they reveal whether the problem is
  localised or widespread. \\[4pt]
3 & Branch-flow / loading error &
  Direct congestion-management target.  Line flows $S_{ij}$ and
  transformer loadings are computed from the predicted state; errors
  here translate directly to wrong congestion decisions. \\[4pt]
4 & $|V|$ RMSE, $\theta$ RMSE &
  Closeness to the NR reference voltage.  Necessary but not sufficient.
  A prediction can be close in MSE but physically inconsistent
  (see Section~\ref{sec:quadrant}). \\
\bottomrule
\end{tabular}
\end{table}

\paragraph{Why $\|\cdot\|_\infty$ before mean.}
AC-PF feasibility is an \emph{all-bus} requirement: every constrained bus
must satisfy \eqref{eq:dp}--\eqref{eq:dq}.  One bus with a large violation
is sufficient to declare the state infeasible, regardless of how well all
other buses behave.  Mean/p95 then answers whether the failure is isolated
or systemic.

\paragraph{Literature grounding.}
\begin{itemize}[noitemsep]
  \item MATPOWER defines AC power flow as solving the mismatch equations
    $g(x)=0$ for bus voltage variables~\cite{matpower}.  The residual being
    small is the definition of a solution, not an optional diagnostic.
  \item OPF formulations enforce nodal power balance alongside line-flow
    and voltage limits~\cite{opfreview}.  Physical feasibility is a hard
    constraint, not a soft metric.
  \item Approximate OPF solutions often violate AC equations and require
    explicit feasibility restoration~\cite{feasrestore}.  This motivates
    reporting power-balance residuals separately from voltage accuracy.
  \item Physics-informed neural AC-OPF papers report worst-case constraint
    violations as primary metrics~\cite{pinnopf}, directly supporting
    Priority 1.
\end{itemize}

%==============================================================================
\section{The Four-Quadrant Classification}
\label{sec:quadrant}
%==============================================================================

Every predicted voltage state falls into one of four quadrants defined by
accuracy (voltage RMSE vs.\ NR reference) and feasibility (PQ residual):

\begin{table}[h]
\centering
\caption{Four-quadrant classification of predicted voltage states.}
\label{tab:quadrant}
\begin{tabular}{p{3.5cm} p{4.5cm} p{4.5cm}}
\toprule
 & \textbf{Low PQ residual} (feasible) & \textbf{High PQ residual} (infeasible) \\
\midrule
\textbf{Close to NR} (accurate) &
  \cellcolor{bestgreen}
  \textbf{Best.} Prediction matches NR and satisfies AC equations.
  Line flows and constraint margins are trustworthy. &
  \cellcolor{warnyellow}
  \textbf{Risky.} Numerically close-looking, but power-balance violated.
  Derived line flows unreliable, especially on stiff buses. \\[8pt]
\textbf{Far from NR} (inaccurate) &
  \cellcolor{neutralblue}
  \textbf{Physically feasible, operationally uncertain.}
  Satisfies AC equations but may correspond to a different operating point.
  Usable if it is the correct load scenario. &
  \cellcolor{badred}
  \textbf{Worst.} Neither accurate nor physically feasible. \\
\bottomrule
\end{tabular}
\end{table}

\paragraph{The subtle case: Close + High residual.}
This happens in stiff multi-voltage grids where small voltage errors produce
large power mismatches via \eqref{eq:stiff}.  The voltage vector looks
acceptable in RMSE, but no NR solver would accept it as converged.
Congestion conclusions drawn from such a state can be wrong because the
inferred branch flows do not correspond to any valid operating point.

\paragraph{For congestion management, the ranking is:}
\begin{enumerate}[noitemsep]
  \item Close to NR + Low residual --- best.
  \item Far from NR + Low residual --- physically feasible; verify operating point.
  \item Close to NR + High residual --- numerically plausible, physically wrong.
  \item Far from NR + High residual --- worst.
\end{enumerate}

%==============================================================================
\section{Experimental Results}
%==============================================================================

\subsection{Do Not Mix Perturbation Regimes}

The LVN logs contain two different data regimes.  They must be reported
separately: the old low-residual Claude run in
\texttt{lvn\_case300\_configs\_share\_grid\_20260531\_141537} used
\textbf{Perturbation A}, while the later LVN experiments use
\textbf{snapshot-envelope} parquets.  The residual scale is therefore not a
pure model effect; it also reflects the sampled operating distribution.

\begin{table}[h]
\centering
\caption{LVN Heo1 log groups and data regimes.  The Perturbation A rows are
  retained as useful historical evidence, but are not directly comparable to
  snapshot-envelope rows.}
\label{tab:regimes}
\setlength{\tabcolsep}{4pt}
\small
\begin{tabular}{p{3.1cm}p{2.9cm}p{4.9cm}p{3.8cm}}
\toprule
Log group & Regime & Parquet evidence & Comparable with \\
\midrule
\texttt{20260531\_141537}
  & Perturbation A, DC start
  & \texttt{LVN\_ppcY\_A\_dc\_compile\_\ldots.parquet}
  & Perturbation A only \\
\texttt{20260531\_141541}
  & Snapshot-envelope, manual flat
  & \texttt{LVN\_snapshot\_envelope\_manual\_flat\_\ldots.parquet}
  & Snapshot-envelope flat-start runs \\
Current 24-job ablation
  & Snapshot-envelope, DC and flat
  & DC and manual-flat snapshot-envelope parquets; residuals evaluated in
    \texttt{complex128}
  & Snapshot-envelope runs only \\
\bottomrule
\end{tabular}
\end{table}

\subsection{Snapshot-Envelope Results}

Table~\ref{tab:results} reports snapshot-envelope results only.  This is the
right table for judging the current LVN snapshot-envelope experiments.

\begin{table}[h]
\centering
\caption{LVN Heo1 snapshot-envelope comparison.  All metrics are validation
  values unless noted otherwise.
  $\Delta P^\infty$, $\Delta Q^\infty$, mean, and p95 are in p.u.\
  (masked: slack excluded from $\Delta P$; slack+PV excluded from $\Delta Q$).
  Baseline rows (grey) are computed directly from the dataset parquet using
  \texttt{complex128} and the same per-unit conversion as training
  ($S_{\rm base}=10^8$\,VA).
  ``---'' means not logged.}
\label{tab:results}
\setlength{\tabcolsep}{3.5pt}
\small
\begin{tabular}{p{3.0cm}rrrrrrrp{2.0cm}}
\toprule
Run & $|V|$ RMSE & $\theta$ & $\Delta P^\infty$ & $\Delta Q^\infty$
    & $\overline{|\Delta P|}$ & $\overline{|\Delta Q|}$
    & p95 $|\Delta P|$ & Notes \\
\midrule
%--- reference baselines ---
\rowcolor{gray!15}
\multicolumn{9}{l}{\textit{Reference baselines (snapshot-envelope, complex128, 36k cases)}} \\
\rowcolor{gray!15}
DC init ($u_{\rm start}$)
  & 4.03e-2 & 0.337\ensuremath{^\circ} & 6.28e2 & 5.01e3 & 2.810 & 11.81 & ---
  & Start of DC-init dataset \\
\rowcolor{gray!15}
NR solution ($u_{\rm newton}$, DC)
  & \textit{0} & \textit{0\ensuremath{^\circ}} & \textit{9.54e-10} & \textit{3.60e-9}
  & \textit{2.46e-12} & \textit{1.04e-11} & ---
  & \textit{Converged reference} \\
\rowcolor{gray!15}
Flat start ($u_{\rm start}$)
  & 4.03e-2 & 3.440\ensuremath{^\circ} & 6.32e2 & 5.01e3 & 2.814 & 11.81 & ---
  & Start of flat dataset \\
\rowcolor{gray!15}
NR solution ($u_{\rm newton}$, flat)
  & \textit{0} & \textit{0\ensuremath{^\circ}} & \textit{9.99e-10} & \textit{3.18e-9}
  & \textit{2.48e-12} & \textit{1.03e-11} & ---
  & \textit{Converged reference} \\
\midrule
%--- model runs ---
\multicolumn{9}{l}{\textit{Old Claude snapshot-envelope manual-flat runs (20260531\_141541)}} \\[2pt]
\rowcolor{neutralblue}
Claude snapshot flat
  \newline\quad K40 dhi16 attn8
  & 7.56e-2 & 3.697\ensuremath{^\circ} & 171.3 & 518.3
  & 1.686 & 1.594 & 2.644
  & Snapshot flat \\[3pt]
\rowcolor{neutralblue}
Claude snapshot flat
  \newline\quad K40 dhi16 attn10
  & 7.55e-2 & 3.691\ensuremath{^\circ} & 169.3 & 518.3
  & 1.691 & 1.580 & 3.052
  & Snapshot flat \\[3pt]
\multicolumn{9}{l}{\textit{Restored old-fixed Armijo, PINN-only A100 reruns (20260614, batch 16)}} \\[2pt]
\rowcolor{bestgreen}
Restored old-fixed
  \newline\quad flat K40 dhi16 attn8
  & 4.02e-2 & 3.879\ensuremath{^\circ} & 169.4 & 518.3
  & 1.686 & 1.595 & 2.947
  & PINN only, no \texttt{vn\_feat} \\[3pt]
\rowcolor{neutralblue}
Restored old-fixed
  \newline\quad flat K40 dhi16 attn10
  & 3.95e-2 & 3.480\ensuremath{^\circ} & 459.1 & 3725
  & 3.510 & 9.760 & 6.660
  & PINN only, no \texttt{vn\_feat} \\[3pt]
\rowcolor{bestgreen}
Restored old-fixed
  \newline\quad flat K40 dhi16 attn10 + \texttt{vn}
  & 3.95e-2 & 3.438\ensuremath{^\circ} & \textbf{41.2} & 519.7
  & \textbf{0.353} & \textbf{1.493} & \textbf{0.903}
  & Best restored flat residual \\[3pt]
\rowcolor{bestgreen}
Restored old-fixed
  \newline\quad DC K40 dhi16 attn8
  & 3.99e-2 & 1.592\ensuremath{^\circ} & 167.3 & 518.2
  & 1.639 & 1.582 & 3.040
  & PINN only, no \texttt{vn\_feat} \\[3pt]
\rowcolor{neutralblue}
Restored old-fixed
  \newline\quad DC K40 dhi16 attn10
  & 4.00e-2 & 1.715\ensuremath{^\circ} & 211.4 & 518.4
  & 1.709 & 2.245 & 3.868
  & PINN only, no \texttt{vn\_feat} \\[3pt]
\rowcolor{neutralblue}
Restored old-fixed
  \newline\quad DC K40 dhi16 attn10 + \texttt{vn}
  & 3.92e-2 & 1.225\ensuremath{^\circ} & 303.4 & 513.9
  & 3.133 & 2.012 & 6.491
  & Lowest restored DC angle \\[3pt]
\midrule
\multicolumn{9}{l}{\textit{Current snapshot-envelope 24-job ablation (complex128 residuals; last validation epoch)}} \\[2pt]
\rowcolor{bestgreen}
moregeom/reject \texttt{dc\_pinn\_mse}
  & \textbf{2.02e-2} & \textbf{0.154\ensuremath{^\circ}} & 902.7 & 1733
  & 11.60 & 27.90 & 38.75
  & Best $|V|,\theta$ current \\[3pt]
\rowcolor{neutralblue}
moregeom/reject \texttt{dc\_pinn\_only}
  & 4.73e-2 & 1.548\ensuremath{^\circ} & 705.8 & 801.9
  & 9.957 & 11.59 & 16.36
  & Best max residual current \\[3pt]
moregeom/reject \texttt{flat\_pinn\_only}
  & 7.28e-2 & 3.502\ensuremath{^\circ} & 1842 & 3841
  & 16.71 & 26.97 & 12.49
  & Best current flat max residual \\[3pt]
\rowcolor{warnyellow}
Student pure MSE
  \newline\quad 260601-5c0d
  & \textbf{2.18e-2} & 0.750\ensuremath{^\circ} & --- & ---
  & --- & --- & ---
  & Best $|V|$; no res.\ logged \\[3pt]
\rowcolor{warnyellow}
Student PINN/mixed
  \newline\quad 260601-8258
  & 2.43e-2 & 2.032\ensuremath{^\circ} & --- & ---
  & --- & --- & ---
  & No res.\ logged \\
\bottomrule
\end{tabular}
\end{table}

\subsection{Perturbation A Results}

Table~\ref{tab:pertA} preserves the old Claude low-residual results, but
labels them as Perturbation A.  These runs used
\texttt{LVN\_ppcY\_A\_dc\_compile\_siNR\_36000\_NR\_branchrows\_directSI.parquet}
from log group \texttt{20260531\_141537}.  They should be used as evidence
that the architecture can obtain low residuals on an easier/narrower
operating distribution, not as a snapshot-envelope benchmark.

\begin{table}[h]
\centering
\caption{Old Claude Perturbation A validation results
  (\texttt{20260531\_141537}).  All residuals are in p.u.}
\label{tab:pertA}
\setlength{\tabcolsep}{4pt}
\small
\begin{tabular}{p{3.2cm}rrrrrrp{3.0cm}}
\toprule
Run & $|V|$ RMSE & $\theta$ & $\Delta P^\infty$ & $\Delta Q^\infty$
    & $\overline{|\Delta P|}$ & $\overline{|\Delta Q|}$ & Notes \\
\midrule
\rowcolor{bestgreen}
Perturb A DC K40 dhi16 attn8
  & 4.504e-2 & 0.229\ensuremath{^\circ} & \textbf{16.65} & 88.58
  & \textbf{0.192} & 0.552
  & Historical low-residual run \\
Perturb A DC K40 dhi16 attn10
  & 4.505e-2 & 0.236\ensuremath{^\circ} & 29.23 & \textbf{70.22}
  & 0.208 & \textbf{0.513}
  & Historical low-residual run \\
Perturb A DC K40 dhi24 attn8
  & 4.484e-2 & 0.239\ensuremath{^\circ} & 17.22 & 73.20
  & 0.226 & 0.561
  & Historical low-residual run \\
Perturb A DC K40 dhi24 attn10
  & 4.507e-2 & 0.256\ensuremath{^\circ} & 41.44 & 74.43
  & 0.304 & 0.608
  & Historical low-residual run \\
\bottomrule
\end{tabular}
\end{table}

\subsection{Reading the Table}

\paragraph{Reference baselines.}
The grey rows are computed directly from the LVN snapshot-envelope parquets
using \texttt{complex128}, the same per-unit conversion as training
($S_{\rm base}=10^8$\,VA, \texttt{share\_grid=True}), and the same masked
residual convention.  They establish three reference levels:

\begin{itemize}[noitemsep]
  \item \textbf{NR solution} ($u_{\rm newton}$): the ground-truth converged
    state. In \texttt{complex128}, $\|\Delta P\|_\infty$ and
    $\|\Delta Q\|_\infty$ are approximately $10^{-9}$\,p.u.  This is the
    target any surrogate should aim for.
  \item \textbf{DC init} ($u_{\rm start}$): the starting point for DC-init
    runs. $\|\Delta P\|_\infty = 628$\,p.u.,
    $\|\Delta Q\|_\infty = 5013$\,p.u.  A model must reduce these
    significantly to be useful.
  \item \textbf{Flat start} ($u_{\rm start}$): the starting point for flat
    runs. $\|\Delta P\|_\infty = 632$\,p.u., $\|\Delta Q\|_\infty = 5013$\,p.u.
    The residual is similar to DC init, but the angle is much worse
    ($3.440^\circ$ vs.\ $0.337^\circ$).
\end{itemize}

\paragraph{Old Claude snapshot-envelope flat runs.}
The old snapshot-envelope flat-start runs in \texttt{20260531\_141541}
achieve $\|\Delta P\|_\infty \approx 170$\,p.u.\ and
$\|\Delta Q\|_\infty \approx 518$\,p.u.  These are the strongest
snapshot-envelope feasibility numbers currently found in the old logs, but
they are manual-flat snapshot-envelope runs, not Perturbation A and not
DC-start snapshot-envelope.

\paragraph{Current 24-job ablation.}
Among the current complex128-evaluated jobs,
\texttt{moregeom/reject + dc\_pinn\_mse} gives the best voltage/angle
accuracy, while \texttt{moregeom/reject + dc\_pinn\_only} gives the best
maximum residual.  Both are snapshot-envelope results, but neither reaches
the lower residual level of the old snapshot-envelope flat run.

\paragraph{Perturbation A low-residual result.}
The old Claude Perturbation A run with K40, dhi16, attn8 achieves
$\|\Delta P\|_\infty = 16.65$\,p.u.  This is the lowest residual found in
the historical logs, but it belongs to a different data regime.  It should
not be cited as the best snapshot-envelope result.

\paragraph{Student pure MSE --- best supervised accuracy, unknown feasibility.}
The student MSE run achieves the best $|V|$ RMSE ($2.18 \times 10^{-2}$)
and a competitive $\theta$ RMSE ($0.75^\circ$), but power-balance residuals
were not logged.  For congestion management this is a risk: the model may be
close to the NR reference in MSE while violating AC equations (the ``Close
+ High residual'' quadrant).  Residual evaluation on saved checkpoints is
required before this run can be recommended for operational use.

\subsection{Run Ranking for Congestion Management}

\paragraph{Snapshot-envelope ranking.}
\begin{enumerate}
  \item \textbf{Old Claude snapshot-envelope flat K40 dhi16 attn8/10}:
    lowest snapshot-envelope residuals found in the logs
    ($\Delta P^\infty \approx 170$, $\Delta Q^\infty \approx 518$).
  \item \textbf{moregeom/reject + dc\_pinn\_only}: best max residual among
    current complex128-evaluated jobs.
  \item \textbf{moregeom/reject + dc\_pinn\_mse}: best voltage/angle accuracy
    among current complex128-evaluated jobs.
  \item \textbf{Student pure MSE}: best supervised accuracy, but feasibility
    unknown until residuals are evaluated.
\end{enumerate}

\paragraph{Perturbation A ranking.}
Within Perturbation A only, the old Claude K40 dhi16 attn8 run is the best
low-residual reference.  It is valuable as a proof that the model family can
learn a physically consistent correction, but it should not be mixed with
snapshot-envelope conclusions.

%==============================================================================
\section{Recommended Reporting Table}
%==============================================================================

\begin{table}[h]
\centering
\caption{Recommended metric set for every GNN AC-PF surrogate evaluation.
  Masking: slack excluded from $\Delta P$; slack + PV excluded from
  $\Delta Q$.}
\label{tab:reporting}
\begin{tabular}{lllc}
\toprule
Metric & Unit & Statistics & Priority \\
\midrule
$|\Delta P|$ (PV+PQ masked) & p.u. / MW   & max, p95, mean & \textbf{1} \\
$|\Delta Q|$ (PQ masked)    & p.u. / MVAr & max, p95, mean & \textbf{1} \\
Convergence rate
  & \% buses $< 10^{-3}$\,p.u. & --- & 2 \\
Branch-flow error
  & p.u. / MVA & MAE, max & 3 \\
$|V|$ error & p.u. & MAE, max & 4 \\
$\theta$ error & deg & MAE, max & 4 \\
NR polish steps & count & mean over test set & (optional) \\
\bottomrule
\end{tabular}
\end{table}

\paragraph{NR polish step count (optional).}
Starting from the surrogate prediction, run pandapower to convergence and
count NR iterations.  Compare against a flat start ($|V|=1$, $\theta=0$).
This is a practical bridge metric: a surrogate that halves NR iterations is
concretely useful even if its raw residual is not at solver tolerance.

\paragraph{Always report both max and mean/p95.}
$\|\cdot\|_\infty$ is the feasibility alarm: one bad bus fails the whole
state.  Mean and p95 characterise the distribution: is the failure
localised or widespread?  Both are needed.  Reporting only mean can mask
a single catastrophically bad bus; reporting only max can be dominated by
one edge-case outlier.

%==============================================================================
%==============================================================================
\section{Numerical Precision: Why Residual Evaluation Requires
  \texttt{complex128} on LVN}
%==============================================================================

\subsection{The Catastrophic Cancellation Problem}

The AC power-flow residual is defined as:

\begin{equation}
  \Delta S_i = S_i^{\rm set} - V_i \cdot \overline{(YV)_i}
  \label{eq:residual}
\end{equation}

At a converged NR solution the two terms in \eqref{eq:residual} are
nearly equal, so $\Delta S_i \approx 0$.  Computing a tiny difference
between two large, nearly equal numbers is the classical
\emph{catastrophic cancellation} scenario in floating-point arithmetic.

The severity of the cancellation depends on the scale of the intermediate
quantities.  In LVN, admittance entries reach $|Y_{ij}| \sim 10^6$\,p.u.\
(see Section~\ref{sec:stiff}).  The matrix-vector product $YV$ therefore
involves values of order $10^6$, and the true residual is of order $10^{-9}$
--- a ratio of $10^{15}$.

\paragraph{\texttt{complex64} precision budget.}
\texttt{complex64} (IEEE 754 single-precision) stores approximately
$7$ significant decimal digits.  For a quantity of magnitude $10^6$,
the representable resolution is:

\begin{equation}
  \delta_{\rm float32} \approx 10^6 \times 10^{-7} = 0.1
\end{equation}

Rounding errors of order $0.1$\,p.u.\ accumulate in the matrix-vector
product, completely swamping a true residual of $10^{-9}$\,p.u.

\paragraph{\texttt{complex128} precision budget.}
\texttt{complex128} (IEEE 754 double-precision) stores approximately
$15$--$16$ significant decimal digits.  For the same $10^6$-scale
quantities:

\begin{equation}
  \delta_{\rm float64} \approx 10^6 \times 10^{-16} = 10^{-10}
\end{equation}

This is sufficient to resolve residuals of order $10^{-9}$\,p.u., which
is the true NR convergence level.

\paragraph{Simple analogy.}
Suppose the true injected and received powers at a bus are:
\begin{align*}
  S^{\rm set} &= 1\,000\,000.123\,456\;\text{p.u.}\\
  V \cdot \overline{(YV)} &= 1\,000\,000.123\,455\;\text{p.u.}\\
  \text{True residual} &= 0.000\,001\;\text{p.u.}
\end{align*}
In \texttt{float32}, both values round to $1\,000\,000$ and the residual
becomes $0$ (or $\sim 0.1$ from inconsistent rounding).  In
\texttt{float64}, both are stored to 15 digits and the difference is
recovered correctly.

\subsection{Observed Effect on LVN Baseline Residuals}

Table~\ref{tab:c128} shows the baseline state residuals computed in
\texttt{complex128} from the LVN snapshot-envelope parquets, using the
same per-unit convention as training ($S_{\rm base}=10^8$\,VA).

\begin{table}[h]
\centering
\caption{LVN Heo1 residuals: \texttt{complex128} baselines (top) and
  24-job ablation runs (bottom, last validation epoch before
  timeout/failure).  Groups: \texttt{default}/\texttt{moregeom}
  $\times$ \texttt{geometric\_safe}/\texttt{reject}.
  All residuals in p.u.\ (masked convention;
  $S_{\rm base}=10^8$\,VA).}
\label{tab:c128}
\setlength{\tabcolsep}{2.5pt}
\footnotesize
\begin{tabular}{llrrrrrrrr}
\toprule
Group / State & Case & ep
  & $|V|$ RMSE & $\theta$
  & $\Delta P^\infty$ & $\Delta Q^\infty$
  & $\overline{|\Delta P|}$ & $\overline{|\Delta Q|}$
  & p95 $|\Delta P|$ \\
\midrule
\multicolumn{10}{l}{\textit{Reference baselines (complex128, 36k cases)}} \\
DC init ($u_{\rm start}$)       & --- & --- & 4.03e-2 & 0.337\ensuremath{^\circ} & 6.28e2 & 5.01e3 & 2.810 & 11.81 & --- \\
\rowcolor{bestgreen}
NR solution (DC)                & --- & --- & \textit{0} & \textit{0\ensuremath{^\circ}} & \textit{9.54e-10} & \textit{3.60e-9} & \textit{2.46e-12} & \textit{1.04e-11} & --- \\
Flat start ($u_{\rm start}$)    & --- & --- & 4.03e-2 & 3.440\ensuremath{^\circ} & 6.32e2 & 5.01e3 & 2.814 & 11.81 & --- \\
\rowcolor{bestgreen}
NR solution (flat)              & --- & --- & \textit{0} & \textit{0\ensuremath{^\circ}} & \textit{9.99e-10} & \textit{3.18e-9} & \textit{2.48e-12} & \textit{1.03e-11} & --- \\
\midrule
\multicolumn{10}{l}{\textit{default / geometric\_safe}} \\
& dc\_pinn\_only  & 10 & 4.717e-2 & 1.014\ensuremath{^\circ} & 2.806e5 & 1.568e5 & 2.782e3 & 1.266e3 & 1.921e3 \\
& dc\_mse\_only   & 10 & 3.248e-2 & 0.528\ensuremath{^\circ} & 1.585e5 & 9.241e4 & 1.309e3 & 7.105e2 & 3.853e2 \\
& dc\_pinn\_mse   & 12 & 4.611e-2 & 0.937\ensuremath{^\circ} & 2.696e5 & 1.422e5 & 2.555e3 & 1.072e3 & 1.710e3 \\
& flat\_pinn\_only& 12 & 7.961e-2 & 3.773\ensuremath{^\circ} & 2.176e5 & 1.037e5 & 2.435e3 & 1.050e3 & 2.123e3 \\
& flat\_mse\_only & 10 & 3.764e-2 & 1.270\ensuremath{^\circ} & 1.084e5 & 3.095e4 & 7.332e2 & 3.111e2 & 8.271e1 \\
& flat\_pinn\_mse & 12 & 7.908e-2 & 3.741\ensuremath{^\circ} & 2.129e5 & 5.163e4 & 1.837e3 & 4.442e2 & 1.473e3 \\
\midrule
\multicolumn{10}{l}{\textit{moregeom / geometric\_safe}} \\
& dc\_pinn\_only  & 11 & 4.145e-2 & 0.553\ensuremath{^\circ} & 3.155e3 & 3.695e3 & 32.47 & 29.21 & 8.984 \\
& dc\_mse\_only   & 10 & 3.672e-2 & 0.184\ensuremath{^\circ} & 3.515e3 & 4.727e3 & 18.79 & 17.46 & 2.957 \\
& dc\_pinn\_mse   & 11 & 4.053e-2 & 0.339\ensuremath{^\circ} & 3.903e3 & 4.620e3 & 33.10 & 25.57 & 1.939 \\
& flat\_pinn\_only& 12 & 7.133e-2 & 3.395\ensuremath{^\circ} & 3.685e3 & 4.254e3 & 41.22 & 21.51 & 21.86 \\
& flat\_mse\_only & 10 & 6.030e-2 & 2.737\ensuremath{^\circ} & 6.500e3 & 4.862e3 & 42.17 & 32.51 & 5.650 \\
& flat\_pinn\_mse & 13 & 7.187e-2 & 3.422\ensuremath{^\circ} & 2.860e3 & 4.305e3 & 32.66 & 19.29 & 2.070 \\
\midrule
\multicolumn{10}{l}{\textit{default / reject}} \\
& dc\_pinn\_only  & 30 & 4.121e-2 & 0.522\ensuremath{^\circ} & 4.280e3 & 4.154e3 & 17.39 & 29.63 & 4.070 \\
& dc\_mse\_only   & 30 & 4.071e-2 & 0.333\ensuremath{^\circ} & 6.283e2 & 5.013e3 & 2.810 & 11.81 & 5.960e-2 \\
& dc\_pinn\_mse   &  1 & 4.155e-2 & 0.610\ensuremath{^\circ} & 4.642e3 & 4.774e3 & 35.60 & 29.21 & 8.229 \\
& flat\_pinn\_only& 30 & 7.190e-2 & 3.414\ensuremath{^\circ} & 4.955e3 & 4.071e3 & 19.97 & 32.16 & 4.607 \\
& flat\_mse\_only & 30 & 7.229e-2 & 3.438\ensuremath{^\circ} & 6.315e2 & 5.013e3 & 2.814 & 11.81 & 1.586e-1 \\
& flat\_pinn\_mse &  1 & 7.157e-2 & 3.395\ensuremath{^\circ} & 4.899e3 & 4.701e3 & 15.10 & 15.86 & 1.988 \\
\midrule
\multicolumn{10}{l}{\textit{moregeom / reject}} \\
\rowcolor{bestgreen}
& dc\_pinn\_mse   & 14 & \textbf{2.020e-2} & \textbf{0.154\ensuremath{^\circ}} & 9.027e2 & 1.733e3 & 11.60 & 27.90 & 38.75 \\
\rowcolor{neutralblue}
& dc\_pinn\_only  & 29 & 4.731e-2 & 1.548\ensuremath{^\circ} & \textbf{7.058e2} & \textbf{8.019e2} & 9.957 & 11.59 & 16.36 \\
& dc\_mse\_only   & 15 & 3.958e-2 & 0.199\ensuremath{^\circ} & 3.286e3 & 4.938e3 & 28.22 & 20.47 & 0.987 \\
& flat\_pinn\_only& 29 & 7.281e-2 & 3.502\ensuremath{^\circ} & 1.842e3 & 3.841e3 & 16.71 & 26.97 & 12.49 \\
& flat\_mse\_only & 11 & 6.064e-2 & 2.728\ensuremath{^\circ} & 3.286e3 & 4.795e3 & 7.124 & 11.92 & 0.326 \\
& flat\_pinn\_mse & 12 & 3.558e-2 & 1.110\ensuremath{^\circ} & 2.962e3 & 1.920e3 & 33.68 & 22.36 & 35.39 \\
\bottomrule
\end{tabular}
\end{table}

\subsection{Key Observations}

\paragraph{NR residuals are at machine precision.}
The NR solution residuals ($\Delta P^\infty \approx 10^{-9}$\,p.u.,
mean $|\Delta P| \approx 10^{-12}$\,p.u.) confirm that $u_{\rm newton}$
is genuinely converged.  These numbers are only recoverable in
\texttt{complex128}; in \texttt{complex64} they would appear as $\sim0.1$
due to cancellation error --- an artefact, not physics.

\paragraph{DC init vs.\ flat start: angle matters, residual does not.}
DC initialisation gives a far better angle estimate ($0.337^\circ$ vs.\
$3.440^\circ$), but the AC residuals are almost identical:
\begin{align*}
  \text{DC init:}   &\quad \text{mean}\,|\Delta P| = 2.810,\quad \|\Delta Q\|_\infty = 5013 \\
  \text{Flat start:}&\quad \text{mean}\,|\Delta P| = 2.814,\quad \|\Delta Q\|_\infty = 5013
\end{align*}
DC initialisation improves the voltage/angle starting point but does not
make the state AC-feasible.  Both starts are equally far from satisfying
\eqref{eq:residual}; the model must do the real work of reducing residuals.

\paragraph{Practical implication for model evaluation.}
The previous evaluation used \texttt{complex64}, which yielded
NR-solution residuals of $\sim 10^{-1}$\,p.u.\ --- many orders of
magnitude above truth.  The current snapshot-envelope ablation table uses
\texttt{complex128} residual evaluation.  Older historical log entries are
kept with their original logged metrics and are therefore best used for
within-regime comparison, not as exact numerical baselines against the
\texttt{complex128} 24-job ablation.

\subsection{Recommended Precision Policy}

\begin{itemize}[noitemsep]
  \item \textbf{Training}: \texttt{complex64} is acceptable.  The physics
    loss is a sum of squared residuals optimised by gradient descent;
    absolute residual magnitude matters less than its gradient direction.
  \item \textbf{Residual evaluation on LVN}: always use \texttt{complex128}.
    The true NR residual is $\sim10^{-9}$\,p.u.; \texttt{complex64} cannot
    resolve below $\sim10^{-1}$\,p.u.\ on this grid.
  \item \textbf{HV/MV grids}: \texttt{complex64} is sufficient.
    Admittances are smaller ($|Y_{ij}| \ll 10^6$), so cancellation is
    less severe and the NR residual ($\sim10^{-8}$\,p.u.) is above the
    \texttt{float32} noise floor.
\end{itemize}

\subsection{24-Job Ablation: Key Observations}

None of the 24 jobs reached final test evaluation (all ended in TIMEOUT or
FAILED); metrics above are from the last validation epoch before termination.

\paragraph{Best overall balance --- \texttt{moregeom/reject + dc\_pinn\_mse} (green).}
Lowest Val RMSE ($2.020 \times 10^{-2}$) and best angle error ($0.154^\circ$)
across all 24 jobs.  $\|\Delta P\|_\infty = 9.027 \times 10^2$\,p.u.\ and
$\|\Delta Q\|_\infty = 1.733 \times 10^3$\,p.u.\ are roughly $100\times$
below the \texttt{default/geometric\_safe} runs.

\paragraph{Best max residual --- \texttt{moregeom/reject + dc\_pinn\_only} (blue).}
Lowest worst-case residuals ($\|\Delta P\|_\infty = 7.058 \times 10^2$,
$\|\Delta Q\|_\infty = 8.019 \times 10^2$), but angle accuracy is worse
($1.548^\circ$ vs.\ $0.154^\circ$).  Preferred when a single-bus outlier
is the dominant concern; \texttt{dc\_pinn\_mse} is better for general accuracy.

\paragraph{\texttt{default/reject} MSE-only cases --- frozen at start.}
\texttt{dc\_mse\_only} and \texttt{flat\_mse\_only} in \texttt{default/reject}
show $\|\Delta P\|_\infty \approx 6.3 \times 10^2$ and mean $|\Delta P|
\approx 2.81$, matching the flat-start initialisation residual (baselines
above).  These runs did not reduce AC residuals despite 30 epochs, indicating
training collapse or dataset mismatch --- a pattern absent in
\texttt{moregeom/reject}, suggesting \texttt{moregeom} perturbation diversity
is critical for the MSE-only variant.

%==============================================================================
\section{Summary}
%==============================================================================

\begin{itemize}[noitemsep]
  \item AC-PF feasibility is defined by power-balance equality constraints
    \eqref{eq:dp}--\eqref{eq:dq}, not by closeness to a reference voltage.
  \item In stiff multi-voltage grids, small voltage errors can produce large
    power mismatches via $|\Delta S| \approx |Y_{ii}| \cdot \epsilon$.
  \item For congestion management, PQ residuals ($\|\Delta P\|_\infty$,
    $\|\Delta Q\|_\infty$, mean, p95) are the primary metrics; voltage
    RMSE is secondary.
  \item Among \textbf{snapshot-envelope} runs, the old Claude manual-flat
    K40 dhi16 attn8/10 logs have the lowest historical residuals found so
    far; among the current complex128 24-job ablation, \texttt{moregeom/reject}
    is the strongest family.
  \item The very low-residual old Claude K40 dhi16 attn8 result belongs to
    \textbf{Perturbation A}.  It is useful historical evidence, but must not
    be mixed with snapshot-envelope conclusions.
  \item The student MSE run has the best supervised accuracy but lacks
    residual evaluation; it should not be used for operational purposes
    without residual verification.
  \item The recommended evaluation table (Table~\ref{tab:reporting}) reports
    residuals first, convergence rate second, branch flows third, and
    voltage/angle RMSE last.
  \item \textbf{Residual evaluation on LVN must use \texttt{complex128}.}
    In \texttt{complex64}, cancellation errors of $\sim10^{-1}$\,p.u.\
    swamp the true NR residual of $\sim10^{-9}$\,p.u., making model
    residuals and the NR baseline numerically incomparable.
    \texttt{complex128} baseline residuals are provided in
    Table~\ref{tab:c128}; GNN model results will be added once evaluated.
\end{itemize}

\begin{thebibliography}{9}

\bibitem{matpower}
R.~D. Zimmerman, C.~E. Murillo-S\'{a}nchez, and R.~J. Thomas,
``MATPOWER: Steady-state operations, planning, and analysis tools for
power systems research and education,''
\emph{IEEE Trans.\ Power Syst.}, vol.~26, no.~1, pp.~12--19, 2011.

\bibitem{opfreview}
D.~Molzahn and I.~Hiskens,
``A survey of relaxations and approximations of the power flow equations,''
\emph{Foundations and Trends in Electric Energy Systems}, vol.~4,
no.~1--2, pp.~1--221, 2019.

\bibitem{feasrestore}
S.~Pineda and J.~M. Morales,
``Feasibility of candidate solutions in gas and power networks,''
\emph{IEEE Trans.\ Power Syst.}, 2022.
(See also: arXiv:2209.04399, ``Restoring AC Power Flow Feasibility
from Relaxed and Approximated OPF Models.'')

\bibitem{pinnopf}
A.~S. Zamzam and N.~D. Sidiropoulos,
``Physics-informed machine learning for power flow with
applications to DC-optimal power flow,''
\emph{arXiv:2110.02672}, 2021.

\end{thebibliography}

\end{document}
